\section{Arbitrary roots on a collar}\label{sec:arbitrary-roots}

The first stage used roots of order $\ell$, a prime different from $p$.
To treat arbitrary ramification, we now construct an auxiliary branched
cover on which a coordinate has a root of any prescribed positive order.
Its own horizontal ramification indices still divide $\ell$.
The construction uses a curve whose stable skeleton has parallel edges
and multiplication on its Jacobian. Throughout this section,
$v=-\log_p|\cdot|$.

We first isolate the step that must work on every sheet. For an analytic
map with a specified lift to a uniformizing torus, pullback by
multiplication by $m$ extracts $m$-th roots of the lift's character
values. The period lattice controls how local roots glue.

\begin{lemma}[Roots after division on a totally degenerate abelian variety]\label{lem:division-roots}
Let $J$ be an abelian variety over $\Omega$ with analytic uniformization
$q:\mathbb T^{\rm an}\to J^{\rm an}$, where $\mathbb T$ is a split torus
and the kernel $\Lambda\subset\mathbb T(\Omega)$ is a free lattice acting by
translation. Suppose $q$ is a topological covering and a local analytic
isomorphism. Let $A$ be an analytic space, let $a:A\to\mathbb T^{\rm an}$
be analytic, and let $m\ge1$ be an integer. Set
\[
 B=A\times_{q\circ a,J^{\rm an},[m]}J^{\rm an}.
\]
For every character $\chi$ of $\mathbb T$, the pullback of $\chi\circ a$
to $B$ has an analytic $m$-th root.
\end{lemma}

\begin{proof}
Let $\pi:B\to A$ and $b:B\to J^{\rm an}$ be the projections, and
write $a_B=a\circ\pi$. Locally on $B$, lift $b$ to an analytic map
$t$ to $\mathbb T^{\rm an}$. The identity $[m]b=q\circ a_B$ says that
$t^m/a_B$ takes values in $\Lambda$. This is a locally
constant function because $\Lambda$ is discrete. Its image in $\Lambda/\Lambda^m$ is
independent of the local lift: replacing $t$ by $t\rho$ changes $t^m/a_B$
by $\rho^m$, for $\rho\in\Lambda$. We therefore obtain an analytic-domain
decomposition into finitely many clopen subsets $B_{\bar\lambda}$,
indexed by $\Lambda/\Lambda^m$.

Choose a representative $\lambda\in\Lambda$ for each coset. On
$B_{\bar\lambda}$, adjust each local lift by a locally constant period
so that $t^m=\lambda a_B$. Two adjusted lifts on an overlap differ by a
period $\rho$ with $\rho^m=1$. The lattice $\Lambda$ is torsion free, so
$\rho=1$. Thus the adjusted lifts glue to a single analytic map
$t_{\bar\lambda}:B_{\bar\lambda}\to\mathbb T^{\rm an}$. Choose
$d_\lambda\in\Omega^\times$ with $d_\lambda^m=\chi(\lambda)$.
The function $\chi(t_{\bar\lambda})/d_\lambda$ is an $m$-th root of
$\chi(a_B)$. These functions glue across the disjoint clopen subsets.
No assumption that $m$ is prime to $p$ is used.
\end{proof}

To obtain a root of an annular coordinate from the preceding lemma,
we arrange that the pullback of a torus character has slope one.
It then differs from the coordinate by a constant and an analytic unit
close to one.  Division on the Jacobian supplies a root of the character
on the whole inverse annulus; a sufficiently long annulus leaves room
for a middle collar on which the remaining unit also admits a root.

\begin{proposition}[An auxiliary branched cover]\label{prop:arbitrary-roots}
Let $m\ge1$ be an integer, choose an integer
\[
 D>v_p(m)+\frac1{p-1},
\]
and let $\alpha\in\overline{k}^{\times}$ satisfy $a=v(\alpha)>2D$.
There is a smooth, proper, connected curve $Q/\overline{k}$ and a finite
map $Q\to\mathbb P^1_{\overline{k}}$ whose ramification indices divide
$\ell$, such that the pullback of the coordinate $T$ has an analytic
$m$-th root on the entire inverse image, over $\Omega$, of
\[
 A^\circ=\{D<v(T)<a-D\}.
\]
\end{proposition}

\begin{proof}
Let $Z$ be the smooth projective normalization of
\[
 y^\ell=\frac{T(T-1)}{T-\alpha}.
\]
The four distinct points $0,\alpha,1,\infty$ have orders $1,-1,1,-1$
in the divisor of the displayed rational function. The cyclic cover
is connected, has degree $\ell$, and is totally ramified at precisely
these points. Riemann--Hurwitz gives
\[
 2g(Z)-2=-2\ell+4(\ell-1),\qquad g(Z)=\ell-1.
\]

Here is a split semistable model of $Z$. Work first over the integers
of a sufficiently large finite extension of $k$ containing $\alpha$.
The blowup of the model of the line along $(T,\alpha)$ has two rational
components and one node, with local equation $TU=\alpha$. The inner
component has coordinate $S=T/\alpha$; the four branch sections are
disjoint and smooth. The Kummer function reduces to $T-1$ on the outer
component and to $-S/(S-1)$ on the inner component. Each of these
degree-$\ell$ covers of a rational component has just two totally
ramified points and is rational. Near the node the Kummer function is
\[
 \frac{T-1}{1-U},
\]
a unit whose reduction at the node is $-1$. Since $\ell$ is invertible,
normalization is \'etale there, with $\ell$ split copies of the node
after enlarging the finite constant field. Away from the node and
branch sections the same unit criterion gives \'etaleness. At a branch
section the function or its inverse is a simple relative parameter,
so the normalization is smooth there. Thus there are exactly two
rational components, joined by $\ell$ nodes, each of thickness
$v(\alpha)=a$. They are stable because $\ell\ge3$.

Consequently the skeleton $\Gamma$ consists of two vertices and
$\ell$ parallel edges $e_1,\ldots,e_\ell$, each of length $a$.
The node annuli $A_j$ are disjoint copies of
$\{|\alpha|<|T|<1\}$ with coordinate $T$. This last assertion can also
be checked directly: on that annulus both $T$ and $\alpha/T$ have
absolute value less than one, and the two factors in
$-(1-T)/(1-\alpha/T)$ have convergent $\ell$-th roots. Hence the
inverse image is exactly $\ell$ copies of the annulus.

We next apply nonarchimedean uniformization of abelian varieties,
developed by Raynaud and Bosch--L\"utkebohmert
\cite{RaynaudUniformization,BoschLutkebohmert,BoschUniformization}.
Let $J$ be the Jacobian of $Z$, and choose an Abel--Jacobi map
$i:Z\to J$ using a point over $\overline{k}$. Total degeneration gives
a toric analytic uniformization
\[
 q:\mathbb T^{\rm an}\longrightarrow J_\Omega^{\rm an},
 \qquad \ker(q)=\Lambda\subset\mathbb T(\Omega).
\]
Its character lattice is $M=H_1(\Gamma,\mathbb Z)$; its tropical target
is $\operatorname{Hom}(M,\mathbb R)$. The period lattice is free, and
$q$ is a topological covering and local analytic isomorphism.
These are the Jacobian uniformization conventions of
\cite[Sections~4.1--4.2 and~4.5.2]{BakerRabinoff}. The compatibility of the classical
and tropical Abel--Jacobi maps \cite[Proposition~6.1 and Corollary~6.6]{BakerRabinoff} identifies the latter with integration
on $\Gamma$ using the edge-length pairing.

The open Berkovich annulus $A_j$ strongly deformation retracts to its
open skeletal interval \cite[Section~2.3]{BakerPayneRabinoff}
and is locally path connected. It is thus
simply connected, so $i|_{A_j}$ lifts to an analytic map
$\widehat i_j:A_j\to\mathbb T^{\rm an}$: first lift continuously along
the covering $q$, then use its analytic local inverses. Orient all
edges by increasing $v(T)$. For $j'\ne j$, the cycle
$c_j=e_j-e_{j'}$ is an element of $M$. If $\chi_j$ is its character,
then $u_j=\chi_j\circ\widehat i_j$ is an analytic unit. The derivative
of $v(u_j)$ along $e_j$, with respect to $v(T)$, is exactly one: the
integral of the harmonic form represented by $c_j$ over a segment of
$e_j$ is the length of that segment. Lifts of the tropical map differ
only by a constant period, which does not change this derivative.

By Lemma~\ref{lem:annular-unit}, an analytic unit on an open annulus
has a unique dominant Laurent monomial at every interior Gauss
radius; its exponent is constant and is the slope of its log-norm. It follows that
\[
 u_j=c_j' T(1+h_j),\qquad
 h_j=\sum_{n\ne0}h_{j,n}T^n,
 \qquad c_j'\in\Omega^\times.
\]
Here $c_j'T$ is the actual coefficient term of degree one.
Dominance at every radius $r\in(|\alpha|,1)$ gives
$|h_{j,n}|r^n<1$. Taking limits at the endpoints yields, for $n>0$,
\[
 |h_{j,n}|\le1,\qquad |h_{j,-n}|\le|\alpha|^n.
\]
Writing $s=v(T)$, the ultrametric inequality consequently gives
\[
 |h_j|\le p^{-\min\{s,a-s\}}<p^{-D}
 \qquad\text{on }A_j\cap\{D<s<a-D\}.
\]
Figure~\ref{fig:parallel-edges} displays the cycle that supplies the
slope-one character and the trimming that gives this bound.

\begin{figure}[hbp]
\centering
\begin{tikzpicture}[
  x=1cm,y=1cm,
  every node/.style={font=\small},
  directed/.style={postaction={decorate},decoration={markings,
    mark=at position 0.56 with {\arrow{Stealth}}}},
  selected/.style={blue!65!black,line width=1pt,directed},
  ordinary/.style={black!60,line width=.7pt,directed}
]
  \node[font=\small\bfseries] at (2,2.0) {Parallel edges of the skeleton};
  \draw[selected] (0,0) .. controls (1,1.5) and (3,1.5) .. (4,0);
  \draw[selected] (0,0) -- (4,0);
  \draw[ordinary] (0,0) .. controls (1,-1.5) and (3,-1.5) .. (4,0);
  \fill (0,0) circle (2pt);
  \fill (4,0) circle (2pt);
  \node[above] at (2,1.13) {$e_j$};
  \node[above] at (2,0) {$e_{j'}$};
  \node at (2,-.60) {$\vdots$};
  \node[below] at (2,-1.13) {other edges};
  \node[left,align=right] at (0,0) {$v(T)=0$};
  \node[right,align=left] at (4,0) {$v(T)=a$};
  \node[align=center] at (2,-2.00)
    {$\ell$ edges, each of length $a$\\cycle $c_j=e_j-e_{j'}$};

  \node[font=\small\bfseries] at (9.4,2.0) {The middle collar on one edge};
  \draw[-{Stealth},black!60] (6.6,1.05) -- (12.2,1.05)
    node[midway,above] {increasing $v(T)$};
  \draw[black!25,line width=2pt] (6.6,0) -- (12.2,0);
  \draw[blue!65!black,line width=2pt] (7.9,0) -- (10.9,0);
  \draw (6.6,-.10) -- (6.6,.10);
  \draw (12.2,-.10) -- (12.2,.10);
  \draw[blue!65!black,fill=white] (7.9,0) circle (2.5pt);
  \draw[blue!65!black,fill=white] (10.9,0) circle (2.5pt);
  \node[below] at (6.6,-.12) {$0$};
  \node[below] at (7.9,-.12) {$D$};
  \node[below] at (10.9,-.12) {$a-D$};
  \node[below] at (12.2,-.12) {$a$};
  \node[above,text=blue!65!black] at (9.4,.12) {$A^\circ$};
  \node[align=center] at (9.4,-1.22)
    {$u_j=c_j'T(1+h_j)$\\$|h_j|<p^{-D}$ on $A^\circ$};
  \node at (9.4,-2.00) {$D>v_p(m)+\dfrac{1}{p-1}$};
\end{tikzpicture}
\caption{The construction in Proposition~\ref{prop:arbitrary-roots},
shown schematically. All edges are oriented by increasing $v(T)$;
the cycle $e_j-e_{j'}$ gives a character unit $u_j$ of slope one
on the annulus corresponding to $e_j$. Division by $m$ on the Jacobian
supplies a root of $u_j$ on the whole inverse annulus. Trimming its
skeletal interval to $D<v(T)<a-D$ makes $h_j$ small enough for
$1+h_j$ to have an $m$-th root, and hence gives a root of $T$.}
\label{fig:parallel-edges}
\end{figure}


On every closed subannulus this estimate has a
strict uniform margin. The logarithm and exponential therefore
converge locally uniformly and
\[
 b_j=\exp\bigl(m^{-1}\log(1+h_j)\bigr)
\]
is analytic there, with $b_j^m=1+h_j$. Indeed, since
$D>1/(p-1)$, the logarithm satisfies
$|\log(1+h_j)|\le |h_j|$, and the valuation of the exponent is
greater than $D-v_p(m)>1/(p-1)$.

For all $\ell$ annuli at once, form the same algebraic finite \'etale cover
\[
 Z'=Z\times_{i,J,[m]}J\longrightarrow Z.
\]
It is finite \'etale for every positive integer $m$, since the ground
field has characteristic zero. Lemma~\ref{lem:division-roots} applied
to $a=\widehat i_j$ gives an analytic $m$-th root of $u_j$ on the
whole inverse image of $A_j$. Dividing it by $b_j$ and an $m$-th root
of $c_j'$ gives an analytic $m$-th root of $T$ above the middle
collar. The annuli $A_j$ are disjoint, so this gives a root on its
entire inverse image in $Z'$. Choose any connected component $Q$ of
$Z'$. Its map to $Z$ is finite \'etale and surjective, so it is a smooth
proper connected curve with the required property. Its composite
with $Z\to\mathbb P^1$ has only ramification indices $1$ or $\ell$.
All algebraic maps used to define $Q$ are over $\overline{k}$; the
analytic roots and torus lifts are only required over $\Omega$.
\end{proof}
